\subsection{Even terminals whose shortest residual arm is at least four}\label{proof:R2-C022}
Let $4\le s_1<t_1<\cdots<s_j<t_j$ be any finite even residual list,
$j\ge1$. Fix $F=\prod_{i=1}^jh_{s_i}f_{t_i}$, all $L\ge1$
connector edges toward the receiving root, and its arm $a\ge3$
outward. For every independent choice of these integers,
\begin{equation}\tag{R22.1.1}\label{eq:R22:claim}
\E[f_a(T^LF)^4]\le1-d/100,
\qquad \E[E_of_a\{T^L(E_tF)\}^4]\le1-d/100.
\end{equation}
Use the finite receiving bounds of R2-C007,
the covariance and Taylor identities
\eqref{eq:R17:covariance}--\eqref{eq:R17:taylor}, and
$\kappa=19/32$. The favorable full-domain bounds are
\[
\operatorname{Cov}(C_a,q)\le-\kappa d+d^2<0,\quad
\operatorname{Cov}(C_a,\bar f_2)\le-3d/16+\theta d^2\le-5d/96.
\]
For the second, $\operatorname{Re}(G_aG_2)\ge1-2x\ge3/16$ and
the actual scalar correction is at most $d\delta_2$.
The true $G_2$ norm is $\sqrt{1+x}$, rather than one.
For $C_a+Y+m\le18/5$, the continuous cubic inequality
\eqref{eq:R12:cubic} gives curvature $2453/250<10$ if $m\le1$,
and $162189/16000<41/4$ if $m\le21/20$.
The low-range exponential payments use
\eqref{eq:R14:exponential}; all have $d\le1/16$ and $Ad/2<1$,
so the deficit exceeds $(A-B)d/3$.

\eqref{eq:L96:four}, in the actual stated direction, gives
\[
\begin{gathered}
h_4f_t\le e^{R_{4,t}-69d/20},\quad h_sf_t\le e^{R_{s,t}-16d/5}\ (s\ge5),\\
R_{s,t}=\sum_{r=s+1}^tK^r\one.
\end{gathered}
\]
Thus $F\le e^{R-\Gamma}$ for one complete finite disjoint-interval
mask $R$, $\Gamma\ge16jd/5$; retain $\Gamma\ge69d/20$ for
$j=1,s_1=4$.
For a finite zero--one polynomial, separating the positive and negative
subsum in each real and imaginary coordinate and taking the larger
absolute value gives
\[
|Q(i\sqrt x)|^2\le\frac{1+x}{(1-x^2)^2}
\le32768/16245<101/50.
\]
The infinite geometric sum only bounds the actual finite mask.
Writing $R=BV_0$, its smallest original power is five, so
\begin{equation}\tag{R22.3.3}\label{eq:R22:masknorms}
\begin{gathered}
\|V_0\|_2^2\le(101/50)\theta^3d<d/7,\quad \sigma:=\E R^2\le(101/50)\theta^4d<d/18,\\
|R|^2<d/7\le1/21<1/16.
\end{gathered}
\end{equation}
Hence $|R|<1/4$. The positive Taylor-tail ratio is at most $1/12$,
and the integral remainder treats negative $R$, giving
$e^R\le1+R+(5/8)R^2$.
The original signed moments $D_i=\E B^iR$ obey
\begin{equation}\tag{R22.3.4}\label{eq:R22:signed}
D_0\in[-d/5,d/12],\quad D_1\in[-d/12,d/5],\quad
D_i\le d/12\ (i\ge2).
\end{equation}
Indeed $\E K^{2r}\one=(-1)^r\delta_r$ and odd moments vanish.
Same-sign moment indices are spaced by two, giving the sharper bounds
$D_0\in[-C_\beta\delta_3,C_\beta\delta_4]$,
$D_1\in[-C_\beta\delta_4,C_\beta\delta_3]$,
$D_i\le C_\beta\delta_4$ for $i\ge2$, where $C_\beta=1024/855$.
Here $C_\beta\theta^2=169/855<1/5$ and
$C_\beta\theta^3=2197/27360<1/12$.
The actual first row and propagation are
\begin{equation}\tag{R22.3.6}\label{eq:R22:propagation}
\begin{gathered}
T(e^R)\le q+BR+D_0+(5/4)\sigma\le q+BR+D,
\quad D=11d/72,\\
e^\Gamma T^LF\le Y_L=\bar f_L+B^LR+D\bar f_{L-1}
+\sum_{i=1}^{L-1}D_i\bar f_{L-1-i}\ge0.
\end{gathered}
\end{equation}
Individual signed summands need not be nonnegative.
The finite convolution coefficients satisfy
$\gamma_r=\beta_r-D_{r-1}\ge d/2$ for $r\ge3$ because
$\kappa-1/12=49/96>1/2$; indices one and two will be paid separately.
From \eqref{eq:R22:masknorms} and $\|B\|^2\le\theta$,
\begin{equation}\tag{R22.3.8}\label{eq:R22:tails}
\begin{array}{c|rrrrrr}
i&1&2&3&4&5&6\\\hline
\|B^iR\|_2/\sqrt d\text{ upper bound}
&5/32&1/10&1/16&1/25&1/40&1/60.
\end{array}
\end{equation}
These are rational-square comparisons with $\theta^i/18$; all later
finite powers only reduce the bound.

On $0<d\le1/16$, first let $L=1$ and $Y=q+BR+D$.
Its actual mean satisfies
$1+5d/72\le m\le1+127d/360\le21/20$,
and $\operatorname{Var}Y\le(1369/1024)d$.
The covariance is at most
$-15d/32+(9/8)(5/32)d=-75d/256$.
Since $m\ge1$, the negative Taylor term is at most $-75d/64$.
The row bound $Y\le q+(6/25)\sqrt d+D$ gives
$C_a+Y+m\le3+(36/25)\sqrt d+(91/180)d\le18/5$;
the increment is at most $9/25+91/2880<3/5$.
Using curvature $41/4$, the exponential constants are
\[
\begin{gathered}
A=18965/1440=64/5+57/32-4(127/360),\\
B=51329/4096=(41/4)(1369/1024)-75/64,\quad C=64/5.
\end{gathered}
\]
$C\ge A/2$, $A-B=117715/184320>3/5$, and $Ad/2<1$,
so the deficit exceeds $d/5$.

For $L=2$, use the actual
$Y=\bar f_2+B^2R+Dq+D_1$.
Its mean satisfies $m\le1-203d/360$.
Since $|D_2|=|\langle B^2\one,R\rangle|\le(5/32)d$ and
$D_1\ge-d/12$, it also satisfies
$m\ge1-(313/288)d\ge9/10$.
The actual random vector is $(G_2+D)(B)u+B^2R$;
$\|G_2\|\le\sqrt{45/32}<19/16$, $D\le11/1152<1/100$ and
\eqref{eq:R22:tails} give $\operatorname{Var}Y\le(169/100)d$.
The $Dq$ covariance is nonpositive, and
\[
\begin{gathered}
\operatorname{Cov}(C_a,\bar f_2)\le-(83/512)d,\quad |\operatorname{Cov}(C_a,B^2R)|\le(17/16)(1/10)d,\\
\operatorname{Cov}(C_a,Y)\le-(143/2560)d.
\end{gathered}
\]
Therefore $4m^3\operatorname{Cov}(C_a,Y)\le-3d/20$, since
$4(9/10)^3(143/2560)=416988/2560000>3/20$.
The row bound $Y\le(1+D)q+(37/32)\sqrt d+d/5$ gives
$C_a+Y+m\le3-(7/120)d+(377/160)\sqrt d
\le3+377/640<18/5$.
With curvature ten, the payment constants are
\[
A=4849/288=64/5+57/32+4(203/360),\quad
B=67/4=10(169/100)-3/20,\quad C=64/5.
\]
$A-B=25/288$, $C\ge A/2$, $Ad/2<1$, giving deficit
$>25d/864>d/100$.

For $L=3,4,5$ replace signed coefficients only before nonnegative
actual path factors, obtaining the nonnegative upper bounds
\begin{equation}\tag{R22.6.1}\label{eq:R22:short}
\begin{aligned}
\overline Y_3&=\bar f_3+B^3R+D\bar f_2+(d/5)q+d/12,\\
\overline Y_4&=\bar f_4+B^4R+D\bar f_3+(d/5)\bar f_2+(d/12)q+d/12,\\
\overline Y_5&=\bar f_5+B^5R+D\bar f_4+(d/5)\bar f_3
+(d/12)\bar f_2+(d/12)q+d/12.
\end{aligned}
\end{equation}
Their means are the actual means after these replacements.
Use the correct short expansions already displayed in Module R2-C021;
in this range $\E\bar f_4\le1-83d/32$,
$\|G_4-d\|\le1$ and $\|P_5\|\le1-2d$, where
$P_5=G_5-d(2+z)$ and
$|P_5|^2-(1-2d)^2=x[-1+2d+d^2+(1-2d)x-x^2+x^3]\le0$.
The signed mean bounds and whole-vector norm bounds give
\[
\begin{array}{c|c|c}
L&\E\overline Y_L\text{ upper bound}
&\text{random-vector norm}/\sqrt d\text{ upper bound}\\\hline
3&1-533d/360&1+(23/60)d+1/16<11/10\\
4&1-19d/10&1+(857/1800)d+1/25<11/10\\
5&1-2d&1+(193/360)d+1/40<11/10.
\end{array}
\]
The mean losses are respectively
$2-11/72-1/5-1/6=533/360$,
$83/32-11/72-1/5-1/4=2867/1440>19/10$, and
$11/4-11/72-1/5-1/3=743/360>2$.
The vector estimates retain the true $G_2$ norm bound $6/5$.
Their actual pointwise bounds are
\[
\begin{aligned}
\overline Y_3&\le(1+127d/360)q+(21/16)\sqrt d-11d/12,\\
\overline Y_4&\le(1-203d/360)q+(11/10)\sqrt d-11d/12,\\
\overline Y_5&\le(1-533d/360)q+(11/10)\sqrt d-\tau+d/12.
\end{aligned}
\]
For the last two square-root coefficients use
$17/16+(23/60)d<11/10$ and
$1+1/25+(857/1800)d<11/10$.
For $L=3$, $C_a+\overline Y_3+m
\le3-(203/120)d+(201/80)\sqrt d<18/5$:
the increment increases for $t=\sqrt d\le1/4$ and has endpoint
$1003/1920<3/5$. For $L=4,5$ discard favorable negative terms to
obtain $3+(23/10)\sqrt d\le3+23/40<18/5$.
All three rows therefore have curvature ten and the common bounds
\[
\begin{gathered}
m\le1-7d/5,\quad \operatorname{Var}\overline Y_L\le121d/100,\\
\operatorname{Cov}(C_a,\overline Y_L)\le(6/5)d.
\end{gathered}
\]
The covariance upper bound is nonnegative, so $m^3\le1$ is legal.
Payment uses $A=20$, $B=169/10=10(121/100)+24/5$, $C=64/5$;
the actual mean exponent $64/5+57/32+28/5$ exceeds twenty.
$A-B=31/10$, $C\ge A/2$, $Ad/2<1$, so the deficit exceeds $31d/30$.

For every finite $L\ge6$, indices one and two in
\eqref{eq:R22:propagation} cannot be absorbed into $\gamma_r$.
Pay their actual positive cost explicitly:
\[
D\bar f_{L-1}+D_1\bar f_{L-2}\le E_0=(127/180)d,
\]
using $\bar f_n\le2$. All remaining $\gamma_r\ge h=d/2$.
Retain the four last actual paths, whose earliest original index is
$L-3\ge3$:
\begin{equation}\tag{R22.7.2}\label{eq:R22:long}
\widehat Y=1-h(4-d)+[G_L-h(3+2z+z^2)](B)u+B^LR+E_0.
\end{equation}
The actual mean coefficient $\gamma_{L+1}\ge h$ is also retained, so
$m\le1-5h+3hd+E_0\le1-8d/5$; its linear loss is at least
$323/180-3/32=2449/1440>8/5$.
The full finite polynomial certificates
\eqref{eq:R18:shiftfour}--\eqref{eq:R18:shifttwo} give random-vector
norm at most $(1-3h+1/60)\sqrt d<(11/10)\sqrt d$.
The pointwise form is
\[
\begin{gathered}
\widehat Y=(1-3h)q+BV-h(1-d)+E_0,\\
V=[G_{L-1}-h(2+z)](B)u+B^{L-1}R,\quad
\|V\|_2\le(1-2h+1/40)\sqrt d.
\end{gathered}
\]
Since $dp+3hq\ge2d$ and $-2d-h(1-d)+E_0\le0$,
$C_a+\widehat Y+m\le3+(89/40)\sqrt d\le3+89/160<18/5$.
Thus the preceding common variance, covariance and payment bounds
also handle every finite $L\ge6$.

It remains to treat $d\ge1/16$ without changing $F$.
For a single pair with $s_1=4$, its original interval and
$\|G_m\|\le6/5$ give
$|R|\le(6/5)\theta^{3/2}\sqrt d<(8/25)\sqrt d\le32d/25$,
so $\Gamma-|R|>2d$ using $\Gamma\ge69d/20$.
For $s_1\ge5$, $|R|\le(6/5)\theta^2\sqrt d<\sqrt d/5\le4d/5$
gives the same conclusion. For $j\ge2$,
$|R|<\sqrt{d/7}\le4d/\sqrt7<8d/5$ and $\Gamma\ge32d/5$.
Consequently the same actual directions satisfy
\begin{equation}\tag{R22.8.1}\label{eq:R22:high}
0\le F\le e^{-2d},\qquad T^LF\le e^{-2d}\bar f_L.
\end{equation}
For $L\ge3$, the precise bare union of Modules R2-C012, R2-C014
and R2-C017 gives $\E[f_a\bar f_L^4]\le1-3d/40$.
The target is at most $e^{-8d}\le1-8d/(1+8d)
\le1-24d/11<1-d/100$.

For $L=1$, on the full $1/16\le d\le1/3$, use
$C_a\le p+3/4$, $q+1=3-p$.
The cubic inequality with $D=15/4,m=1$ gives curvature
$173/16<11$. Since covariance with $q$ is nonpositive,
$\E[f_aq^4]\le1-57d/32+11d=1+295d/32$.
But $e^{8d}\ge1+8d+32d^2$, whose difference from this bound is
at least $d(32d-39/32)\ge25d/32$.
After multiplying by $e^{-8d}$ the deficit is greater than
$25d/864>d/100$, because $8d\le8/3<3$ and $e^3<27$.

For $L=2$, first take $1/16\le d\le1/8$.
The actual mean and variance are $m_2=1-d$,
$\operatorname{Var}\bar f_2=d+\delta_2-d^2\le d(45/32-d)$,
and covariance with $C_a$ is nonpositive.
Since $\bar f_2=q+Bu$ and $\sqrt d-d\le1/4$,
$\bar f_2+m_2\le13/4-p$, $C_a\le p+9/20$, $m_2\le15/16$.
The full pointwise coefficient is at most
\[
\Phi(p)=(p+9/20)[(13/4-p)^2+225/128]\le10.
\]
Indeed $\Phi''=6p-121/10\le-3$ on $[0,3/2]$,
$\Phi(1)=25317/2560<99/10$, $|\Phi'(1)|=189/640<3/10$,
so the tangent increment is below $3/200$.
After $3/2$, $\Phi'(3/2)=-1283/640<0$ and
$\Phi''\le-1/10$, so the remainder decreases.
Taylor expansion gives
\begin{align*}
\E[f_a\bar f_2^4]&\le P(d):=(1-57d/32)(1-d)^4+10d(45/32-d)\\
&=1+(265/32)d+(25/8)d^2-(235/16)d^3
+(65/8)d^4-(57/32)d^5\\
&\le1+(265/32)d+(25/8)d^2.
\end{align*}
The omitted tail is nonpositive for $d\le1/8$.
Thus $e^{8d}-P(d)\ge d[(231/8)d-9/32]\ge195d/128$;
the resulting deficit exceeds $(65/128)d>d/100$ since $e^{-8d}>1/3$.

Finally take $L=2$, $1/8\le d\le1/3$.
Now $C_a\le p+3/4$, $m_2\le7/8$ and the coefficient is bounded by
\[
\Psi(p)=(p+3/4)[(13/4-p)^2+49/32]\le12.
\]
On $[0,1]$, $\Psi''\le-5$,
$\Psi(3/4)=747/64$, $\Psi'(3/4)=9/32$, so the tangent increment
is at most $81/10240$ and the value below twelve.
On $[1,2]$, the convex quadratic derivative has negative endpoint
values $-41/32,-121/32$, so $\Psi$ decreases.
The target is bounded by
\[
e^{-8d}[(1-57d/32)(1-d)^4+12d(45/32-d)].
\]
The first term is at most $(3/8)(4/5)(7/8)^4=7203/40960$,
using $e\ge8/3$. The function
$h(d)=d(45/32-d)e^{-8d}$ decreases here: after removal of the
positive exponential its derivative is
$45/32-2d-(45/4)d+8d^2$, whose derivative is at most $-95/12$
and whose value at $1/8$ is $-1/8$.
Thus $h(d)\le41/(256e)\le123/2048<1/16$.
The total is at most $7203/40960+3/4=37923/40960<14/15\le1-d/5$.
All high rows join the low range at $d=1/16$.
At $d=0$, original leaf integration gives equality, without requiring
$K=0$. Positivity restores all original equal pairs and subdivision
vertices. Full reversal reverses receiving arm, terminal arms,
connectors, and decorations together.
